\begin{frame}{공통 원리: 작업에 대해 ``상대적''}
  \begin{block}{속도와 정밀도는 \textit{작업}에 대해 상대적이다}
    ``이 작업에 \textbf{필요한} 정밀도''와 ``이 작업에
    \textbf{필요한} 경험 수''를 \textbf{먼저 재고}, 그 뒤에
    ``이 환경의 steps/s''를 \textbf{대입}해야 한다.
    \textbf{역순}(환경이 빠르니까/정밀하니까 쓰자)으로 시작하면
    앞 표의 실수를 한다.
  \end{block}
  \vskip0.7em
  \begin{itemize}
    \item \kb{정밀도 = 필요충분 조건} + \kb{속도 = 그 안에서
      비용 절감}
  \end{itemize}
  \vskip0.5em
  {\footnotesize 확인 문제 예시: 13.2 사족보행 1,000만 스텝.
  (a) 단일 CPU MuJoCo(Ant, 약 6,200 steps/s) 환경 스텝만 실
  시간 $\approx$ \textbf{27분}. (b) 정책 계산·업데이트를 환경
  스텝의 2배로 가정 $\Rightarrow$ 총 $\approx$ \textbf{81분}.
  (c) Isaac Sim이 경험률을 2,000배 $\Rightarrow$ (b)의 총 실
  시간은 $\approx$ \textbf{0.04분(약 2.5초)}으로 줄어듬.
  }
\end{frame}
